\ExplSyntaxOn

\int_new:N \l__binp_int
\int_set:Nn \l__binp_int {4}
\tl_new:N \l__bina_tl
\tl_new:N \l__binb_tl
\tl_new:N \l__binhw_tl
\int_new:N \l__bina_int
\int_new:N \l__binb_int
\int_new:N \l__binc_int
\int_new:N \l__bind_int
\fp_new:N \l__bina_fp

\DeclareDocumentCommand \BinaryPrecision {m}
{
  \int_set:Nn \l__binp_int {#1}
}

\cs_new_nopar:Npn \int_to_bin:Nn #1#2
{
  \tl_clear:N #1
  \int_set:Nn \l__bina_int {#2}
  \prg_replicate:nn {\l__binp_int}
  {
    \int_set:Nn \l__binb_int {\int_mod:nn {\l__bina_int} {2}}
    \tl_put_left:Nx #1 {\int_use:N \l__binb_int}
    \int_set:Nn \l__bina_int {\int_div_truncate:nn {\l__bina_int} {2}}
  }
}

\cs_new_nopar:Npn \bin_hamming_weight:NN #1#2
{
  \tl_set_eq:NN \l__binhw_tl #2
  \tl_replace_all:Nnn \l__binhw_tl {0} {}
  \int_set:Nn #1 {\tl_count:N \l__binhw_tl}
}

\cs_new_nopar:Npn \bin_flip_one:NNn #1#2#3
{
  \tl_set_eq:NN #2#1
  \prg_replicate:nn {#3 - 1}
  {
    \tl_replace_once:Nnn #2 {1} {a}
  }
  \tl_replace_once:Nnn #2 {1} {0}
  \tl_replace_all:Nnn #2 {a} {1}
}

\DeclareDocumentCommand \HammingWeight {m m}
{
  \int_to_bin:Nn \l__binhw_tl {#2}
  \bin_hamming_weight:NN #1 \l__binhw_tl
}

\DeclareDocumentCommand \HyperCube {m}
{
  \int_set:Nn \l__binp_int {#1}
  \int_zero_new:c {l__hyper_0_int}
  \int_set:cn {l__hyper_0_int} {-1}
  \int_step_inline:nnnn {1} {1} {#1}
  {
    \int_zero_new:c {l__hyper_##1_int}
    \int_set:cn  {l__hyper_##1_int} {\int_use:c {l__hyper_\int_eval:n {##1 - 1}_int} * (#1 - ##1 + 1) / ##1}
  }
  \fp_set:Nn \l__bina_fp {2^{#1}-1}
  \int_set:Nn \l__binc_int {\fp_to_int:N \l__bina_fp}
  \int_step_inline:nnnn {0} {1} {\l__binc_int}
  {
    \int_to_bin:Nn \l__bina_tl {##1}
    \bin_hamming_weight:NN \l__binb_int \l__bina_tl
    \node[every~ hypercube~ label/.try] (hg- \l__bina_tl) at (\int_use:c {l__hyper_\int_use:N \l__binb_int _int},\int_use:N \l__binb_int) {\texttt{\l__bina_tl}};
    \int_incr:c {l__hyper_\int_use:N \l__binb_int _int}
    \int_incr:c {l__hyper_\int_use:N \l__binb_int _int}
  }
  \int_step_inline:nnnn {0} {1} {\l__binc_int}
    {
      \int_to_bin:Nn \l__bina_tl {##1}
      \bin_hamming_weight:NN \l__binb_int \l__bina_tl
      \int_step_inline:nnnn {1} {1} {\l__binp_int}
      {
        \bin_flip_one:NNn \l__bina_tl \l__binb_tl {####1}
        \tl_if_eq:NNTF \l__bina_tl \l__binb_tl
          { % do nothing if nodes are the same
          }
          {
            \draw[every~ hypercube~ edge/.try] (hg- \l__bina_tl) -- (hg- \l__binb_tl);
          }
      }
    }
  }

\ExplSyntaxOff

% Macro to cross out a single node
\newcommand{\CrossOutNode}[1]{%
  \draw[red, thick]
    ($(#1) + (-0.3, 0.3)$)
    --
    ($(#1) + (0.3, -0.3)$);
  \draw[red, thick]
    ($(#1) + (-0.3, -0.3)$)
    --
    ($(#1) + (0.3, 0.3)$);
}

% Macro to cross out a list of nodes
\newcommand{\CrossOutNodes}[1]{%
  \foreach \n in {#1} {%
    \CrossOutNode{\n}%
  }
}
% Macro to draw a circle around a single node
\newcommand{\CircleNode}[2][RoyalBlue]{%
  \draw[#1, thick]
    ($(#2)$)
    circle (0.5);
}

% Macro to circle a list of nodes
\newcommand{\CircleNodes}[2][RoyalBlue]{%
  \foreach \n in {#2} {%
    \CircleNode[#1]{\n}%
  }
}

% Überwachungstabelle der Aufgabe mit 10 Nachrichtenbits und 4 Kontrollbits
% (Spalten 1 bis 14, Zeilen c1 bis c4); Massstab im tikzpicture setzen
\newcommand{\ueberwachungZehnVier}{%
  \fill[blue!25] (0,0) rectangle (15,1);
  \fill[blue!15] (0,-4) rectangle (1,0);
  \fill[gray!8] (1,-4) rectangle (15,0);
  \foreach \i in {1,...,10} {
      \node at (\i+.5,.5) {\(a_{\i}\)};
    }
  \foreach \i in {1,...,4} {
      \pgfmathtruncatemacro{\col}{10+\i}
      \node[text=orange!80!black] at (\col+.5,.5) {\(c_{\i}\)};
      \node[text=orange!80!black] at (.5,.5-\i) {\(c_{\i}\)};
    }
  % Überwachungsbereiche von c1 bis c4, eine Liste pro Zeile
  \foreach \cols [count=\z] in {{1,2,3,5,6,7,11},{1,2,4,5,8,9,12},{1,3,4,6,8,10,13},{2,3,4,7,9,10,14}} {
      \foreach \col in \cols {
          \node at (\col+.5,.5-\z) {\(\times\)};
        }
    }
  % Trennung Nachrichtenbits / Kontrollbits
  \draw[blue!70!black,very thick] (11,-4) -- (11,1);
  \draw[step=1,black!30,thin] (0,-4) grid (15,1);
  \draw[black!65] (0,-4) rectangle (15,1);
}

% Bitwerte des Codeworts über den Spaltentiteln
% #1: 14 Bits als Komma-Liste, #2: Spalte magenta hervorheben (0 = keine)
\newcommand{\ueberwachungWerte}[2]{%
  \node[font=\small\itshape,anchor=east] at (1,1.5) {Wert};
  \foreach \w [count=\col] in {#1} {
      \ifnum\col=#2
        \node[magenta,font=\small\bfseries\ttfamily] at (\col+.5,1.5) {\w};
      \else
        \node[font=\small\ttfamily] at (\col+.5,1.5) {\w};
      \fi
    }
}

% Anzahl Einsen unter den Kreuzen jeder Zeile rechts der Tabelle,
% ungerade Anzahlen (Kontrollbit meldet) magenta
\newcommand{\ueberwachungEinsen}[1]{%
  \node[font=\small\itshape] at (16.4,.5) {Einsen};
  \foreach \n [count=\z] in {#1} {
      \ifodd\n
        \node[magenta] at (16.4,.5-\z) {\textbf{\n}\ \(\boldsymbol{\times}\)};
      \else
        \node[green!50!black] at (16.4,.5-\z) {\n\ \checkmark};
      \fi
    }
}
